\documentclass[11pt,letterpaper]{article}
\usepackage[T1]{fontenc}
\usepackage{lmodern,amsmath,amssymb,amsthm,mathtools}
\usepackage[margin=1in]{geometry}
\usepackage{microtype}
\usepackage[hidelinks]{hyperref}
\hypersetup{pdftitle={Continuum power avoidance for logarithmically syndetic configurations}}
\ifdefined\pdfinfoomitdate\pdfinfoomitdate=1\fi
\ifdefined\pdftrailerid\pdftrailerid{}\fi
\ifdefined\pdfsuppressptexinfo\pdfsuppressptexinfo=15\fi
\usepackage{enumitem}
\setlength{\parskip}{4pt}
\setlength{\parindent}{0pt}
\setlength{\emergencystretch}{2em}
\newtheorem{theorem}{Theorem}
\newtheorem{lemma}{Lemma}
\newtheorem{proposition}{Proposition}
\newtheorem{corollary}{Corollary}
\newcommand{\R}{\mathbb R}
\newcommand{\Z}{\mathbb Z}
\newcommand{\N}{\mathbb N}
\newcommand{\rhoo}{\operatorname{dens}}
\newcommand{\logtwo}{\log_2}
\DeclareMathOperator{\dist}{dist}
\title{Continuum power avoidance for logarithmically syndetic configurations}
\author{}
\date{7 October 2026}
\begin{document}
\begin{center}
{\LARGE Continuum power avoidance for logarithmically syndetic configurations\par}
\vspace{10pt}
7 October 2026
\end{center}
\vspace{4pt}
\begin{abstract}
For any prescribed countable family of positive configurations whose occupied
dyadic logarithmic bins have bounded gaps near zero, we construct one closed,
nowhere-dense, one-periodic set with measure greater than $1-\varepsilon$ in
every unit interval. Every image with a nonzero leading real power and a
positive-power remainder has infinitely many distinct points outside this
set in every sufficiently small tail, simultaneously for all real exponents,
translations, nonzero coefficients, and finite error constants. No regularity
of the image function is assumed. In particular the same set avoids all affine
null geometric progressions for all ratios in $(0,1)$. We give the full routing
proof and an exact sparse-block obstruction to extending this particular
logarithmic-span output-window mechanism to all sources of positive upper
Banach logarithmic density.
\end{abstract}

This is a complete written mathematical proof with an independent AI-assisted
mathematical audit. It is not proof-assistant formalization, human peer review,
a novelty certification, or a practical small-density construction. The exact
finite programs accompanying the note check boundary and arithmetic interfaces;
the infinite theorem is proved below. The inherited routing mechanisms and the
prior finite parameter-arrangement principle are attributed in Section~9.

\section{Statements and quantifiers}
For $A\subset(0,\infty)$ define
\[
 S(A)=\{j\ge1:A\cap(2^{-j-1},2^{-j}]\ne\varnothing\}.
\]
We call $A$ \emph{logarithmically syndetic near zero} if there are integers
$G,J\ge1$ such that $S(A)$ meets $\{j,j+1,\ldots,j+G-1\}$ for every integer
$j\ge J$. This is stronger than positive logarithmic upper Banach density.
There is no measurability or cardinality hypothesis on $A$.

\begin{theorem}\label{thm:main}
Let $(A_\ell)$ be a nonempty prescribed countable family of logarithmically
syndetic configurations. For every $0<\varepsilon<1$ there is a closed,
nowhere-dense, one-periodic $E\subset\R$ such that
\[
 |E\cap[x,x+1]|>1-\varepsilon\qquad(x\in\R),
\]
and the following holds simultaneously for every $\ell$, every real
$s,\alpha>0$, every $y\in\R$, every $c\ne0$, and every finite $M\ge0$.
If $f$ is defined on a sufficiently small tail of $A_\ell$ and
\[
 |f(a)-y-ca^s|\le M a^{s+\alpha}
\]
for every sufficiently small $a\in A_\ell$, then for every sufficiently small
$\varrho>0$ the set
\[
 f(A_\ell\cap(0,\varrho))\setminus E
\]
contains infinitely many distinct points. No regularity or measurability of
$f$ is assumed.
\end{theorem}
The functions, exponents, and translations in this statement are not
enumerated. The countable family of configurations is prescribed first.

\begin{corollary}\label{cor:geometric}
Taking $A=\{2^{-n}:n\ge1\}$ gives one such $E$ excluding every affine null
geometric progression $\{y+cq^n:n\ge1\}$, simultaneously for every
$q\in(0,1)$, $y\in\R$, and $c\ne0$, with infinitely many distinct misses in
every tail.
\end{corollary}
Indeed set $s=-\logtwo q$ and use zero remainders. This corollary is a consequence
of the proof; no literature-priority conclusion is attached to it.

The main local statement is the following compact-continuum lemma. Here and
below the density of a periodic measurable set is its measure on one period.
\begin{lemma}[Uniform robust blocker]\label{lem:blocker}
Fix $A$ as above, $0<s_0<s_1<\infty$, $\alpha_0>0$, $k\in\Z$, integers
$q,h\ge1$, and $0<p<1/12$. There is an open one-periodic $H$ of density at most
$6p$ such that, for every $x\in\R$, $s\in[s_0,s_1]$, $t\in[1,2]$, and every
prescribed error family on $A\cap(0,2^{-h})$ with
\[
 |e(a)|\le q a^{s+\alpha_0},
\]
some $a$ in that tail satisfies
\[
 x+t2^k a^s+e(a)\in H.
\]
The same $H$ works for all these parameters and error families and can be based
on routing windows at an arbitrarily late output logarithmic position $U$.
\end{lemma}
The error family is arbitrary; in particular its members need not be mutually
independent random variables. All probability, entropy, and perturbation costs
needed for the lemma are specified next.

\section{A uniformly spaced input subsequence}
Choose an occupied index $j_1$ beyond $J,4,h$ and put
\[
 m=\lceil3/s_0\rceil+2.
\]
Given $j_n$, choose an occupied $j_{n+1}$ in
$\{j_n+m,\ldots,j_n+m+G-1\}$, possible by syndeticity, and choose $a_n$ in
its dyadic bin. Then, with $z_n=-\logtwo a_n$,
\[
 a_n\downarrow0,\qquad z_n\in[j_n,j_n+1),\qquad
 3/s_0<z_{n+1}-z_n<m+G.
\]
Every $a_n$ lies in $A\cap(0,2^{-h})$. Only this subsequence is used.

Put $B_0=m+G$ and $D=s_1B_0$. For each $s\in[s_0,s_1]$, the output logarithms
\[
 w_n(s)=sz_n-k
\]
strictly increase to infinity, with successive gaps between $3$ and $D$.
Take $U\ge s_1z_1-k$. Every open interval $(u,u+\ell)$, $u\ge U$, contains
at least $\ell/D-1$ of these output logarithms, uniformly in $s$.
To see this, take the last member at most $u$ and the first member at least
$u+\ell$. If $N_\mathrm{int}$ counts the strictly interior members, the interval
is covered by at most $N_\mathrm{int}+1$ consecutive gaps, each at most $D$.
Thus $\ell\le(N_\mathrm{int}+1)D$. This handles equality at either endpoint.
Consequently, for $\ell\ge2D$ and $\eta=1/(2D)$, there are at least
$\eta\ell$ interior members. The lower gap $3$ also gives an upper count
$1+\ell/3\le\ell$ for integer $\ell\ge2$.

The windows are fixed on the \emph{output} axis. Every exponent supplies its
own active input points in these windows, without altering the random set or
its grids. Fixing input windows instead would require separation at $s_1$
while preserving earlier keys at $s_0$, leading to incompatible scales.

\section{The routing template and stable centers}
Fix $p,\eta$ and choose an integer $M\ge2$ with
\[
 p\eta(M-1)\ge12.
\]
Put $\lambda=p\eta(M-1)/2$ and $\kappa=\lambda-4\log2>0$, where $\log$ is
natural. Choose an integer $d\ge1$ with
\[
 (1-2^{1-M})^d<p;
\]
for example any integer $d>2^{M-1}\log(1/p)$ suffices. Let $K$ be the number
of edges in the complete ordered $M$-ary tree of depth $d$. Choose an integer
$g\ge1$ with $K2^{3-g}<p$. These choices do not depend on $U$ or on unexposed
randomness.

List the edges in preorder. Give every edge leaving height $r$ a window of
length $\ell_r$ and leave exactly $g$ unused integer indices between consecutive
windows. Define
\begin{align*}
 \ell_1&=L,&\sigma_1&=ML+(M-1)g,\\
 \ell_r&=g+\sigma_{r-1},&
 \sigma_r&=M(\ell_r+g+\sigma_{r-1})+(M-1)g\quad(r\ge2).
\end{align*}
Set $T=\sigma_d$ and start the template at integer $U$. For an edge $e$ write
its integer window endpoints as $u_e,v_e=u_e+\ell_r-1$, and let $v_e^*$ be the
last endpoint of its edge and child subtree. Then
\[
 v_e^*-u_e+1\le2\ell_r.
\]
For a leaf this span is $\ell_1$; otherwise it equals
$\ell_r+g+\sigma_{r-1}=2\ell_r$. All $\ell_r\ge L$, and the last endpoint of
the complete template is $U+T-1$.

Input $n$ is active on edge $e$ at exponent $s$ exactly when
\[
 u_e<w_n(s)<v_e+1.                                    \tag{A}
\]
Both inequalities are strict. The count just proved applies even at exponents
where an input is deactivated at a window boundary. For $L\ge\lceil2D\rceil$
every edge has at least $\eta\ell_r$ active tests.

At integer endpoint $b$ use the nested periodic grid of $N_b=2^{b+3}$ cells,
with key $J_b(z)=\lfloor N_b\{z\}\rfloor$. Cells are closed on the left and
open on the right. For each nondefault edge $(P,i)$, $i<M$, assign a table of
independent fair bits at its endpoint grid. Assign each terminal leaf an
independent Bernoulli-$p$ table at its incoming-edge grid. All entries in all
tables are mutually independent. Route a real point to the first nondefault
child whose addressed selector bit is one, and otherwise to the default child
$M$. Let $B$ be the points whose terminal bit is one. This is constant on the
finest grid, whose size is
\[
 N=2^{U+T+2},\qquad \mathbb E\rhoo(B)=p.
\]
For the expectation, condition on all selectors: each fixed point's addressed
terminal bit remains Bernoulli-$p$. The probability space is finite.

Call a center $x$ \emph{stable} if, for every noninitial window starting at
$u=b+g+1$ with $b$ the preceding window's endpoint, the interval
\[
 (x,x+2^{-b-g}]
\]
contains no boundary of that predecessor grid. The unstable density is at most
$K2^{3-g}<p$, since each grid has $2^{b+3}$ boundaries per period and the
interval has length $2^{-b-g}$. At a stable center every active offset satisfies
\[
 t2^ka_n^s=t2^{-w_n(s)}\le2^{1-u_e},
\]
and crosses no predecessor boundary. Grid nesting preserves \emph{all}
preceding keys, including ancestor decisions and preceding siblings.

At its own grid the active points and center have pairwise distinct addresses.
Every active offset is greater than $2^{-u_e-\ell_r}$; the cell width is
$2^{-u_e-\ell_r-2}$. Two active output logarithms differ by more than $3$,
so their offset difference is greater than
$(7/8)2^{-u_e-\ell_r}$, again greater than the cell width. Multiplication by
$t\ge1$ preserves the inequalities. When $U\ge4$, all offsets lie in an
interval of length at most $1/8$, so circular distance is ordinary distance.
Distinctness persists on finer grids and holds at every center.

\section{Conditional independence at a fixed continuum parameter}
Fix a stable $x$ and expose its addressed entry in \emph{every} selector table.
Its route is determined. Before conditioning, its probability of never using
a default child is $(1-2^{1-M})^d<p$. On a center-exposure atom with a first
default at a vertex of height $r$, consider every outgoing nondefault edge
$e_i$ and all active indices at a fixed $(s,t)$.

Earlier-key stability makes such points follow the center through the strict
ancestors and reject the preceding siblings at the default vertex. If the
point's own selector bit is one, it enters child $i$. Define its local test
$Q_{i,n}$ to be this own selector multiplied by the terminal bit obtained by
routing from child $i$. Success implies actual membership in $B$.

There are at least $\eta(M-1)\ell_r$ active tests. Their own selector addresses
are distinct, avoid exposed center entries, and are independent fair bits.
Next condition on \emph{all} selectors. The local terminal addresses are
distinct: different child subtrees or leaves use different tables, while points
sharing a leaf retain distinct keys on its finer grid. Thus the terminal bits
are independent Bernoulli-$p$ variables. If $A_j$ are the tested own selectors,
conditional failure is $(1-p)^{\sum A_j}$. Averaging these fair bits gives
exactly
\[
 (1-p/2)^{\text{number of active tests}}
 \le \exp(-\lambda\ell_r).
\]
Auxiliary selector entries may be shared. Active indices may vary with $s$.
Independence is asserted only at a fixed $(s,t)$ after this conditioning;
routing paths themselves are not assumed independent.

\section{Exact finite representation of the continuum}
Keep $x$ and its center-exposure atom fixed, hence also the first-default vertex
and its outgoing windows. For an edge of length $\ell=\ell_r$, candidate indices
can occur only with
\[
 (u_e+k)/s_1\le z_n\le(u_e+\ell+k)/s_0.
\]
Assume $U>|k|$. The gaps of $z_n$ exceed $3/s_0$, so each edge has at most
$\lceil(U+T+|k|)/3\rceil+1$ candidate indices. The number of candidate
edge-index pairs at this vertex is therefore at most
\[
 P=(M-1)\bigl(\lceil(U+T+|k|)/3\rceil+1\bigr).
\]
For each pair record the activation boundaries
\[
 s=(u_e+k)/z_n,\qquad s=(u_e+\ell+k)/z_n
\]
when they lie in $[s_0,s_1]$, and its local finest-grid boundary curves. A local
test is determined by $J_{v_e^*}(x+t2^{k-sz_n})$. Active offsets are positive
and less than $2^{1-u_e}$, so there are at most
\[
 B_\ell=1+2^{2\ell+3}
\]
possible lifted grid boundaries $\beta=b/N_{v_e^*}-x$ in that offset range.
For every such positive $\beta$ the crossing curve is
\[
 t=\beta2^{-k}2^{z_ns}.
\]
Retain its full portion in $[s_0,s_1]\times[1,2]$. Irrelevant retained portions
only increase the count.

Every exponential curve is simple and monotone. Distinct curves intersect at
most once, since equality becomes affine in $s$ after logarithms. For the same
input, different coefficients give no intersection; equal coefficients give
coincident curves, merged with all labels retained. Vertical activation lines
intersect a curve at most once, as do $t=1,2$. Coincident or parallel boundary
lines are treated in the same way. Equivalently the change $v=\logtwo t$ turns
crossing curves into ordinary lines
\[
 v=\logtwo\beta-k+z_ns
\]
in the rectangle $[s_0,s_1]\times[0,1]$. This is an exact change of parameters,
not a numerical approximation.

Include the rectangle's four boundary arcs. There are at most
$C=P(B_\ell+2)+4$ distinct arcs. Pair intersections number at most
$C(C-1)/2$, and endpoints at most $2C$. Splitting at intersections gives at most
$C^2$ open arc pieces; inserting arcs successively gives at most
$1+C(C+1)/2$ open two-dimensional cells. All vertices, open arc pieces, and
open two-dimensional cells therefore admit at most
\[
 20(C+1)^2
\]
representatives, a conservative bound. Degenerate arcs retain their singleton
points separately; the bound still holds by treating at most one singleton
per candidate arc as a vertex.

On each stratum the complete vector of activation indicators and active local
finest-grid keys is constant. Open cells avoid recorded boundaries. On an
exponential boundary the key is the exact integer-boundary key assigned to the
cell on its right; intersections split that arc. On an activation boundary the
pair is inactive by (A). Parameter endpoints and intersections have individual
representatives. Thus all active local tests are constant on a stratum for
\emph{every} realization of the unexposed entries. Representatives are chosen
before those entries are sampled. No sampled exponent grid is used.

If all active local tests fail at any parameter, they fail at a retained
representative. At every representative the uniform active count and Section~4
give failure probability at most $e^{-\lambda\ell}$. The conditional
all-parameter miss probability is consequently at most
\begin{align}
 &20[P(3+2^{2\ell+3})+5]^2e^{-\lambda\ell}\notag\\
 &\hspace{15mm}\le5120M^2(U+T+|k|+2)^2e^{-\kappa\ell}.       \label{eq:entropy}
\end{align}
Indeed $P\le M(U+T+|k|+2)$ and, for $P\ge1$,
$P(3+2^{2\ell+3})+5\le16P2^{2\ell}$.
The absolute-position entropy cost is genuinely polynomial in $U+T+|k|$, and
its local-window exponential rate is $4\log2$ before the failure factor.
It is not treated as an exponent-independent constant. We have taken no union
bound over uncountably many parameters or over every exposure atom with an
extra entropy factor, and use the local final grid rather than the whole
tree's finest grid.

\section{A noncircular window and buffer schedule}
All $M,d,K,g,\eta,\lambda,\kappa$ are fixed. For integer $U$ set
\begin{align*}
 L(U)=\max\Bigl\{g,\lceil2D\rceil,2,
 \Bigl\lceil\frac{4\log(U+2)+\log(46080M^2/p)+1}{\kappa}\Bigr\rceil\Bigr\}.
\end{align*}
Use this $L$ in the template to obtain $T(U)=\sigma_d$. For the fixed tree,
$T$ is affine in $L$, because
\[
 \sigma_r=2M\sigma_{r-1}+(3M-1)g\quad(r\ge2),\qquad
 \sigma_1=ML+(M-1)g.
\]
Thus $T(U)=O(\log U)$, with a finite tree-dependent constant. No practical
size claim is made.

Choose $U$ arbitrarily large subject to
\[
 U\ge4,\quad U>|k|,\quad U\ge s_1z_1-k,\quad T(U)\le U.    \tag{B}
\]
Such integers exist since $T(U)/U\to0$. For every local $\ell\ge L(U)$,
\eqref{eq:entropy} is bounded by
\[
 46080M^2(U+1)^2e^{-\kappa L(U)}
 \le p e^{-1}\frac{(U+1)^2}{(U+2)^4}<p.                    \tag{C}
\]
The no-default event costs less than $p$, so the probability of an all-parameter
miss by the active finite tests is less than $2p$ at each stable center.
Unstable centers have probability at most one, and integrated miss density is
at most $3p$.

If an input is active on any template window, $w_n(s)>U$ and $s\le s_1$ imply
\[
 2^ka_n^s<2^{-U},\qquad z_n>(U+k)/s_1.
\]
Hence
\[
 q a_n^{s+\alpha_0}
 \le q2^{-k}2^{-U}2^{-\alpha_0(U+k)/s_1}.
\]
Set the strictly positive common buffer radius
\[
 r_U=(q2^{-k}+1)2^{-U}2^{-\alpha_0(U+k)/s_1}.
\]
At the finest grid $N=2^{U+T(U)+2}$, the double-buffer cost is
\[
 4Nr_U=2^{T(U)+4}(q2^{-k}+1)2^{-\alpha_0(U+k)/s_1}.        \tag{D}
\]
Since $T(U)=O(\log U)$ and $\alpha_0U/s_1$ is positive and linear, (D) tends
to zero. Increase $U$ until (D) is less than $p$, keeping (B)--(C).
This is noncircular: $L,T$ are explicit logarithmic functions of $U$, and both
the entropy and buffer inequalities hold for every sufficiently large $U$.
The positive lower bound $\alpha_0>0$ is essential here.

\section{Closed missed centers and repair of every center}
On the circle, lifted periodically, put
\[
 B_1=B+(-r_U,r_U),\qquad B_2=B+(-2r_U,2r_U).
\]
For any union of $N$ grid cells, enlargement by radius $r$ increases density
by at most $2Nr$. Therefore $\mathbb E\rhoo(B_2)\le p+4Nr_U<2p$.

Only finitely many indices can be active anywhere for some $s\in[s_0,s_1]$:
their $z_n$ belong to $[(U+k)/s_1,(U+T+k)/s_0]$. Let $I$ be this finite index
set and let $S_n$ consist of the $s$ for which $w_n(s)$ lies in the union of the
\emph{open} template windows. Define $R$ as the centers for which there is
some $(s,t)$ in the compact rectangle such that
\[
 x+t2^ka_n^s\notin B_1\quad\hbox{for every $n$ with $s\in S_n$}.
\]
For each $n\in I$ the relevant condition is
\[
 [s\notin S_n]\quad\hbox{or}\quad[x+t2^ka_n^s\notin B_1].
\]
This is closed: $S_n$ is relatively open, and the second condition is closed
by continuity of the power map and openness of $B_1$. A finite intersection
remains closed. Projecting this closed subset of the compact product of the
circle and the parameter rectangle shows that $R$ is closed. Closed activation
would invalidate this argument; Section~10 gives an exact counterexample.
No measurable choice of stratum representatives as a function of $x$ is needed
for the pointwise probability estimate.

Since $B\subset B_1$, membership in $R$ at a stable center implies failure
of all local tests at the first-default vertex, unless there is no default.
The preceding sections give $\mathbb E\rhoo(R)\le3p$. A finite-table outcome
therefore has
\[
 \rhoo(B_2)+\rhoo(R)\le5p.
\]
By outer regularity choose an open periodic $V\supset R$ with
$\rhoo(V)\le\rhoo(R)+p$, and put $H=B_2\cup V$. It is open, periodic, and
has density at most $6p$.

Fix $x,s,t$ and any allowed error family. If $x\notin R$, at this parameter
an active finite point lies in $B_1$. Its error is at most $r_U$, so its perturbed
point lies in $B_2$: distance strictly less than $r_U$ from $B$ stays strictly
less than $2r_U$ after an error of size at most $r_U$. Both buffers are open;
equality in the error bound loses no boundary point.

If $x\in R$, then $x\in V$, and along the full selected sequence
$t2^ka_n^s+e(a_n)\to0$. Sufficiently late perturbed points therefore lie in $V$.
These indices need not lie in $I$ or be active on the finite template. This
repairs every exceptional center. In fact convergence is uniform in $s,t$,
because for $a_n<1$ the absolute offset is at most
\[
 2^{k+1}a_n^{s_0}+q a_n^{s_0+\alpha_0}.
\]
The errors require no measurability or mutual dependence assumptions and may
be selected after $H$. This completes the proof of Lemma~\ref{lem:blocker}.

\section{One avoiding set for all real powers}
Apply the lemma to every tuple $(\ell,N,j,k,q,h)$ with integers
$N\ge2$, $j,q,h\ge1$, $k\in\Z$, setting
\[
 [s_0,s_1]=[1/N,N],\qquad\alpha_0=1/j.
\]
Choose $0<p_{\rm tuple}<1/12$ with $12\sum p_{\rm tuple}<\varepsilon$.
The tuple family is countable although each blocker handles a continuum of
$s$. Let $\mathcal U$ be the union of all blockers and their reflections and
put $E=\R\setminus\mathcal U$. The density of $\mathcal U$ is less than
$\varepsilon$, so $E$ is closed, one-periodic, and has the required measure
in every unit interval.

Given $s,\alpha>0$, choose $N$ with $s\in[1/N,N]$ and $j$ with $1/j\le\alpha$.
For $a<1$ and integer $q\ge\max(1,M)$,
\[
 M a^{s+\alpha}\le q a^{s+1/j}.
\]
For $c>0$ write $c=2^kt$ with $t\in[1,2]$. For every sufficiently large $h$,
$f$ is defined and its inequality holds on the whole tail below $2^{-h}$.
The corresponding blocker gives a hit at some $a_h\in A_\ell$, $a_h<2^{-h}$.
For $c<0$ apply the positive argument to $-f$ and use reflected blockers.

Since $M a^\alpha\to0$, eventually
\[
 (|c|/2)a^s\le |f(a)-y|\le(3|c|/2)a^s.
\]
The hits approach $y$ without equaling it; a finite set of non-$y$ values
cannot do this. They therefore contain infinitely many distinct values in
 every sufficiently small tail. This proves the full order of quantifiers
in Theorem~\ref{thm:main}. There is no uncountable union of density budgets or
rational approximation of a real exponent.

Affine maps occur with $s=1$ and zero errors. If $E$ contained an interval, a
small affine image of a bounded nonempty tail of some $A_\ell$ would fit inside
it, contradicting the hitting property. Thus $E$ has empty interior and is
nowhere dense because it is closed. In particular $E\cap[0,1]$ is compact,
has measure greater than $1-\varepsilon$, and has the corresponding tail
avoidance property.

\section{Scope and attributed mechanisms}
The proved statement assumes bounded gaps between occupied logarithmic bins.
It covers all real $s>0$ by countably many compact exponent ranges, and within
each blocker $s$ is an uncountable parameter. The active indices move with $s$.
The proof does not substitute a prescribed countable profile list or a
finite-alphabet bounded-cluster switching argument for this continuum step.

The nested-grid random routing and exceptional-center repair are inherited
from the pinned OpenAI geometric-case family 084 and entry 006v1. Their exact
conditional-independence argument is reproved here. Entry 006v2's logarithmic
profile viewpoint motivates the output-axis choice. The combination used here
is exponent-dependent activation in fixed output windows, boundary-complete
parameter stratification, and a logarithmic-window/positive-power-buffer
schedule. Bounded gaps supply uniform counts at every output location; positive
upper Banach logarithmic density alone does not do so.

The finite parameter-arrangement principle has a primary antecedent in
Kolountzakis--Papageorgiou, Theorem~4.1 and its proof \cite{KP}. They record
grid-hyperplane crossings as equations linear in a transformation parameter
and bound polynomially many grid-address regions before a probability union
bound. Our crossing curves become lines under $v=\logtwo t$; the arrangement
principle itself is not asserted to be new. Their theorem concerns a finite
$P$ with at most $N^\zeta$ points and transforms with at least $\alpha(N)$
$N^{-\beta}$-separated image points, where $\alpha(N)/\log N\to\infty$.
It is not used as a premise here. The moving power-activation windows, complete
boundary strata, absolute-position cost, and positive-power buffers are
proved explicitly above. No exhaustive prior-work assessment or novelty
conclusion is made.

There is no extension here to all positive-upper-Banach-density configurations,
arbitrary prescribed slow remainders, arbitrary $C^1$ diffeomorphisms, or flat
smooth germs. The buffer schedule requires $\alpha_0>0$. In particular profiles
such as $2^{-1/a}$ fall outside the nonzero leading-power class. The construction
is existential with enormous finite parameters; no practical generation or
runtime bound is claimed.

The distinct embedding boundary must be attributed precisely.
Feng--Lai--Xiong, Theorem~1.1 \cite{FLX}, gives a global bi-Lipschitz embedding
with $f'(0)=1$ under $\limsup a_{n+N}/a_n<1$. It does \emph{not} state a global
$C^1$-diffeomorphism conclusion. The global $C^1$ endpoint used for comparison
is separately proved in the pinned 006v1 source \cite{M006}. Such an embedding
for a lacunary sequence may have a remainder violating every positive-power
rate treated here. Neither embedding theorem is a premise of the present proof.

A precise dated consequence concerns Burgin--Goldberg--Keleti--MacMahon--Wang,
Question~1 in arXiv:2210.09284v1 (17 October 2022) \cite{BGKMW}. Their
formulation asks whether every positive-measure measurable set contains some
translated, nonconstant null geometric progression. Take the set $E$ above
for $A=\{2^{-n}:n\ge1\}$ and set $F=E\cap[0,1]$. Then $F$ is compact and
$|F|>1-\varepsilon$. For any $a\ne0$, $b\in\R$, and $q\in(0,1)$, put
$s=-\log q/\log2$ and $f(u)=b+au^s$. Its remainder is zero and
$f(2^{-n})=b+aq^n$, so infinitely many points of every tail lie outside $E$,
hence outside $F$. This gives a negative answer to that exact 2022 formulation.
We do not assert that the question remained open in 2026 or claim priority.
The comparison with later literature is incomplete; no broader resolution of
the Erd\H{o}s similarity conjecture is asserted.

The source pins are \texttt{87868bcb65bf0460a90d6bc2efc3324e46481230} for
\texttt{mxym/math} entries 006v1 and 006v2, and
\texttt{adc7f1241b42e322a6451854ab7e4b4c146bf78a} for
\texttt{openai/math} family 084, especially its full window, routing, and scale
sections \cite{O084}. Family 084 states the theorem for each fixed geometric
ratio and permits the set to depend on that ratio. Entry 006v2 treats a
prescribed countable profile family and does not assert all real exponents
simultaneously. These source comparisons do not establish priority.
Exact component URLs, byte hashes, original notices, and licensing limits are
retained in the accompanying provenance files.

\section{Exact boundary controls}
These controls are not the infinite proof. Take $a_1=1/2,a_2=1/4$, all
$s,t\in[1,2]$, $x\in\R$, and the open holes $H=(1/10,1)+\Z$. With a common
independent error radius $r\ge0$, every parameter succeeds when $r<7/160$,
and the certificate fails at $r=7/160$.

For success put $u=2^{-s}\in[1/4,1/2]$. The positive unperturbed difference
$t(u-u^2)$ lies in $[3/16,1/2]$. With two errors the difference lies in
$[3/16-2r,1/2+2r]$. If $r<7/160$ this interval is strictly contained in
$(1/10,9/10)$. Two points outside the open holes are in closed complementary
gaps $[m,m+1/10]$; their positive difference is either at most $1/10$ or at
least $9/10$, a contradiction.

At $r=7/160$, take $s=2,t=1,x=143/160$ and errors $(-7/160,+7/160)$.
The perturbed points are $11/10$ and $1$, both excluded endpoints. At this radius
all $s<2$ or $t>1$ still succeed, so deleting the parameter boundary loses the
failure. This guards the compact closed exponent/coefficient rectangle and
retained point/curve strata without floating powers or sampled grids.

For the closed-activation counterexample, use one input $a=1/2$, $x=0,t=1$,
activation $[1,2]$, and the hole $(2/5,3/5)$. For $s<1$ the input is inactive,
so the condition ``if active, then missed'' holds. At $s=1$ the input becomes
active and $2^{-s}=1/2$ is in the hole, so the condition fails. Its failure set
is not closed. Open activation $(1,2)$ keeps the limiting point inactive and
makes the implication set closed, as Section~7 requires.

The executable controls check these sharp rational endpoints, curve-intersection
inequalities, positive-error buffer costs, and negative replacements with
explicit exceptions that remain active under Python optimization. Independent
controls enumerate 5,322 activation-boundary/open-stratum representatives,
320 preorder templates, 2,352 entropy inequalities, and 4,096 selector
assignments in an actual two-level nested-grid routing instance. The latter
conditions on all center entries and integrates terminal bits exactly:
distinct-terminal failure is $28561/38416$. Deliberate terminal-key coarsening
instead gives $469225/614656$, with collisions in 732 selector outcomes.
It detects the address-independence loss that would invalidate the product
identity. Sharp parameter-boundary, closed-activation, one-buffer, and
sparse-block controls accompany it. None enumerates the full calibrated tree
or formalizes the infinite analytic proof.

\section{A sparse-block obstruction for this mechanism}
Define
\[
 R_m=2^{2^m},\qquad
 S=\bigcup_{m\ge3}\{R_m,R_m+1,\ldots,R_m+m\},\qquad
 A=\{2^{-j}:j\in S\}.
\]
For integers $L\ge1$ define
$F_S(L)=\sup_{r\in\Z}\#(S\cap\{r+1,\ldots,r+L\})$ and
$d^*(S)=\limsup_{L\to\infty}F_S(L)/L$.
Every length $L$ fits in an occupied block, so $F_S(L)=L$ and $d^*(S)=1$.
Nevertheless uniform occupancy fails for every sufficiently late logarithmic-span
template at some exponent in $[1,2]$.

\begin{proposition}\label{prop:sparse}
For every fixed finite $C>0$, every sufficiently large $U$, and every
$0<T\le C\log U$, some $s\in[1,2]$ has no $j\in S$ satisfying
\[
 U<sj<U+T.
\]
Thus this exponent has no active input anywhere in the template.
\end{proposition}
\begin{proof}
An activating $j$ must lie in $[U/2,U+T]$. For large $U$ we have $T\le U/4$.
Any block intersecting this interval has $R_m\le5U/4$, hence
$m=O(\log\log U)$, and $R_m+m\ge U/2$, hence $R_m\ge U/3$ eventually.
There is at most one intersecting block: the next center is
$R_{m+1}=R_m^2\ge U^2/9>5U/4$, and previous centers are separated in the
same way. If there is no intersecting block, every $s$ works. Otherwise every
exponent activating any member of that block lies in
\[
 \left(\frac{U}{R_m+m},\frac{U+T}{R_m}\right).
\]
Its length is
\[
 \frac{T}{R_m}+\frac{Um}{R_m(R_m+m)}
 \le \frac{3T}{U}+\frac{9m}{U},
\]
which tends uniformly to zero for $T\le C\log U$. It is eventually less than
one, so this open interval cannot cover $[1,2]$. Choosing $s$ outside it proves
the claim.
\end{proof}
This counterexample obstructs the \emph{specific uniform output-window,
logarithmic-span mechanism} for general positive-upper-Banach-density sources.
It does not disprove a continuum avoidance theorem for that larger class.
Longer variable spans, different placement principles, or another routing
architecture are not addressed. The proved theorem remains the bounded-gap
statement.

\begin{thebibliography}{9}
\bibitem{O084} OpenAI/math, \emph{The geometric case of the Erd\H{o}s similarity
conjecture}, family 084, 5 October 2026, pinned revision
\texttt{adc7f1241b42e322a6451854ab7e4b4c146bf78a}.
Main, windows, routing, and scale components are included with exact URLs
and hashes in the package.
\bibitem{M006} mxym/math, \emph{Modulus nonlinear similarity}, entry 006,
versions 1 and 2, pinned revision
\texttt{87868bcb65bf0460a90d6bc2efc3324e46481230}.
Reference components and notices are included in the package.
\bibitem{FLX} Feng, Lai, and Xiong, \emph{Erd\H{o}s similarity problem via
bi-Lipschitz embedding}, arXiv:2312.01319v1, Theorem 1.1.
\url{https://arxiv.org/html/2312.01319v1}.
\bibitem{KP} Mihail N. Kolountzakis and Effie Papageorgiou,
\emph{Large sets containing no copies of a given infinite sequence},
Analysis \& PDE \textbf{18} (2025), 93--108,
\href{https://doi.org/10.2140/apde.2025.18.93}{doi:10.2140/apde.2025.18.93}.
Theorem 4.1 and proof in arXiv:2208.02637v2 (6 October 2023),
\url{https://arxiv.org/html/2208.02637v2\#S4}.
\bibitem{BGKMW} Alex Burgin, Samuel Goldberg, Tam\'as Keleti, Connor MacMahon,
and Xianzhi Wang, \emph{Large sets avoiding infinite arithmetic / geometric
progressions}, arXiv:2210.09284v1 (17 October 2022), Question 1.
\url{https://arxiv.org/html/2210.09284v1\#S1}.
\end{thebibliography}
\end{document}
